\originalchapter{静电边值问题}
\originalsection{镜像法}
点电荷 $q$ 位于接地平面上方高度 $h$, 在平面下方对称位置放虚拟电荷 $-q$. 上半空间中的势为
\[
V=\frac q{4\pi\eps}\left[\frac1{\sqrt{s^2+(z-h)^2}}-\frac1{\sqrt{s^2+(z+h)^2}}\right].
\]
它在 $z=0$ 为零, 在上半空间只有真实点电荷的奇点, 且无穷远消失. 唯一性确保这是解. 镜像电荷只用于域内求解, 平面下半空间的镜像场不是导体内部物理场.
\begin{center}
\begin{tikzpicture}[>=Stealth,scale=.8]
\draw[thick] (-3,0)--(3,0);\node[right] at (3,0) {接地导体 $z=0$};
\fill (0,1.7) circle(2pt);\node[right] at (.2,1.7) {$q$};
\fill (0,-1.7) circle(2pt);\node[right] at (.2,-1.7) {$-q$ (虚拟)};
\draw[dashed] (0,-1.5)--(0,1.5);\draw[<->] (-.7,0)--(-.7,1.7);\node[left] at (-.7,.85) {$h$};
\end{tikzpicture}\par
真实求解区域为上半空间; 图中长度为任取比例.
\end{center}
真实电荷受力为 $-q^2\hat{\vct z}/(16\pi\eps h^2)$. 相互作用势能却是 $\frac12qV_\mathrm{image}=-q^2/(16\pi\eps h)$, 其中一半来自感应电荷随真实电荷移动而改变. 对力直接用库仑定律时没有这个一半.

半径 $a$ 的接地球外有点电荷 $q$, 距球心 $d>a$. 球内镜像为 $q'=-qa/d$, 位置距球心 $d'=a^2/d$, 与真实电荷同轴. 对球面任意点, 到两个电荷的距离满足 $R'=aR/d$, 因而两项势抵消. 若球孤立且总电荷为 $Q$, 再在球心加镜像 $Q-q'$, 保持球面等势且总感应电荷为 $Q$.

\originalsection{分离变量与球谐函数}
边界不是可用少量镜像表达的几何时, 应把拉普拉斯方程分解成与边界相容的模态. 例如矩形区域中两侧为零势, 顶边给定 $f(x)$, 取 $V=X(x)Y(y)$ 得 $X''/X=-Y''/Y=-k^2$. 零侧边要求 $X_n=\sin(n\pi x/a)$, 对应 $Y_n$ 是双曲函数. 在底边 $y=0$ 也接地的条带 $0<y<b$ 中,
\[
V(x,y)=\sum_{n=1}^\infty f_n\sin\frac{n\pi x}{a}\frac{\sinh(n\pi y/a)}{\sinh(n\pi b/a)},\quad
f_n=\frac2a\int_0^af(x)\sin\frac{n\pi x}{a}\dd x.
\]
系数由正交性确定. 分段光滑边界的跳跃点须按适当极限解释, 不能要求级数在互相矛盾的角点同时取两个值.

球坐标中 $\nabla^2=r^{-2}\partial_r(r^2\partial_r)+r^{-2}\Delta_\Omega$. 角算子的本征函数为 $\Delta_\Omega Y_{\ell m}=-\ell(\ell+1)Y_{\ell m}$, 其中 $\ell=0,1,2,\ldots$, $-\ell\le m\le\ell$, 采用 $\int Y_{\ell m}^*Y_{\ell'm'}\dd\Omega=\delta_{\ell\ell'}\delta_{mm'}$. 径向方程 $r^2R''+2rR'-\ell(\ell+1)R=0$ 给 $R=r^\ell$ 或 $r^{-\ell-1}$. 原点有界选择前者, 无穷远消失选择后者. 半径 $a$ 的球面给定 $f(\theta,\phi)$ 时, 内部解为
\[
V(r,\theta,\phi)=\sum_{\ell,m}f_{\ell m}(r/a)^\ell Y_{\ell m},\quad
f_{\ell m}=\int Y_{\ell m}^*f\dd\Omega.
\]
最低阶函数为 $Y_{00}=1/\sqrt{4\pi}$、$Y_{10}=\sqrt{3/(4\pi)}\cos\theta$、$Y_{1,\pm1}=\mp\sqrt{3/(8\pi)}\sin\theta e^{\pm i\phi}$. 不同 $\ell$ 的正交来自角拉普拉斯的自伴性: 对球面积分分部, $\int Y^*\Delta_\Omega Z=\int(\Delta_\Omega Y)^*Z$, 球面没有边界, 两个不同本征值之差乘内积为零. 同阶不同 $m$ 则由 $e^{im\phi}$ 的正交性区分.

轴对称时只需勒让德多项式 $P_\ell(\cos\theta)$, 不再有 $\phi$ 依赖. 令 $x=\cos\theta$, 标准归一化 $P_\ell(1)=1$ 给 $\int_{-1}^1P_\ell P_j\dd x=2\delta_{\ell j}/(2\ell+1)$, 所以 $f(x)=\sum a_\ell P_\ell(x)$ 的系数为 $a_\ell=(2\ell+1)\int_{-1}^1f(x)P_\ell(x)\dd x/2$. 例如 $x^2=P_0/3+2P_2/3$, 同一边界在球内延拓成 $V=1/3+(2/3)(r/a)^2P_2(\cos\theta)$, 常数模态不能误删.

\begin{example}
半径 $a$ 的接地导体球置于均匀外电场 $E_0\hat{\vct z}$. 求球外势及表面电荷.
\end{example}
\begin{solution}
远处的势为 $-E_0r\cos\theta$, 是 $\ell=1$ 模态. 加入消失于无穷远的同阶项, 写 $V=(-E_0r+A/r^2)\cos\theta$. 在 $r=a$ 为零给 $A=E_0a^3$. 径向场在球面为 $3E_0\cos\theta$, 故 $\sigma=3\eps E_0\cos\theta$. 它的总电荷为零, 偶极矩为 $4\pi\eps a^3E_0\hat{\vct z}$.
\end{solution}
\originalsection{从角向微分方程理解球谐}
球面上的角拉普拉斯为
\[
 \Delta_\Omega=\frac1{\sin\theta}\partial_\theta(\sin\theta\,\partial_\theta)
 +\frac1{\sin^2\theta}\partial_\phi^2.
\]
令 $Y=\Theta(\theta)\Phi(\phi)$, $\Delta_\Omega Y=-\lambda Y$. $\phi$ 的周期性要求 $\Phi=e^{im\phi}$、$m\in\mathbb Z$. 用 $x=\cos\theta$, 剩下的方程为
\[
 \frac\dd{\dd x}\left[(1-x^2)\frac{\dd\Theta}{\dd x}\right]
 +\left[\lambda-\frac{m^2}{1-x^2}\right]\Theta=0.
\]
要求南北极处的球面函数正则, 可取 $\Theta=(1-x^2)^{|m|/2}$ 乘一个多项式. 把多项式按幂代入, 最高次项要求 $\lambda=\ell(\ell+1)$, 且 $\ell\ge|m|$ 为整数; 其余系数由递推决定. 对 $m=0$ 就是勒让德方程, 取 $P_\ell(1)=1$. 对 $m\ge0$, 关联函数可定义为
$P_\ell^m(x)=(-1)^m(1-x^2)^{m/2}\dd^mP_\ell/\dd x^m$.
结合 $e^{im\phi}$ 并做单位球面归一化就得到本稿的 $Y_{\ell m}$. 因此 $m$ 是绕轴的角变化次数, $\ell$ 是角拉普拉斯本征值的标号; 它们并非随意引入的两个下标.

径向幂次是角变化对空间延拓的约束. 在球内, $r^\ell$ 在原点正则; 在球外, $r^{-\ell-1}$ 对应相同角变化的远场衰减. 例如二极的 $\ell=1$ 模态由内场 $r\cos\theta=z$ 接到外势 $r^{-2}\cos\theta$, 四极的 $\ell=2$ 由二阶调和多项式接到 $r^{-3}P_2$. 这把边值问题和多极展开接在一起, 也帮助判断漏掉某阶后会丢掉怎样的场.

\begin{example}
二维矩形 $0<x<a,0<y<b$, 左右与底边接地, 顶边取常电势 $V_0$. 写出内部电势, 说明角点条件如何理解.
\end{example}
\begin{solution}
先用左右零边界选 $\sin(n\pi x/a)$, 底边零边界选 $\sinh(n\pi y/a)$. 顶边常数的正弦系数为
\[ f_n=\frac{2V_0}{a}\int_0^a\sin(n\pi x/a)\dd x=\frac{2V_0[1-(-1)^n]}{n\pi}. \]
因此仅有奇数阶,
\[
 V=\frac{4V_0}{\pi}\sum_{j=0}^\infty
 \frac{\sin[(2j+1)\pi x/a]}{2j+1}
 \frac{\sinh[(2j+1)\pi y/a]}{\sinh[(2j+1)\pi b/a]}.
\]
内部每一项调和. 顶边开区间内级数收敛到 $V_0$, 侧边为零. 顶角同时被赋零与 $V_0$ 的两种边界值, 所以不要求角点处取同一个连续值; 物理上该理想尖角需要更细的电极连接模型. 这道例题是本稿的教学补充, 不冒充缺失的 Griffiths 3.15.
\end{solution}

\originalsection{带圆孔球壳：由面荷跳跃决定系数}
8.07第7讲实际给出半径$R$球面在北极挖去锥角$\alpha$球帽的例子. 剩余面荷密度仍为$\sigma_0=Q/(4\pi R^2)$, $Q$表示未挖孔时电荷, 实际剩余电荷为$Q(1+\cos\alpha)/2$, 两者不可混同. 令$x=\cos\theta$, $\sigma(x)=\sigma_0$ ($-1<x<\cos\alpha$), 其余零.

球内正则及球外衰减各选一支, 在$r=R$电势连续使它们同用系数$\widetilde A_\ell$:
$V=\sum_{\ell\ge0}\widetilde A_\ell r_<^\ell r_>^{-\ell-1}P_\ell(x)$, $r_<=\min(r,R)$、$r_>=\max(r,R)$. 法向场差逐项给
$\sum(2\ell+1)\widetilde A_\ell P_\ell=R^2\sigma/\eps$. 用$\int_{-1}^1P_\ell P_j\dd x=2\delta_{\ell j}/(2\ell+1)$投影,
\[
\widetilde A_\ell=\frac{R^2}{2\eps}\int_{-1}^1P_\ell(x)\sigma(x)\dd x
=\frac{Q}{8\pi\eps}\int_{-1}^{\cos\alpha}P_\ell(x)\dd x.
\]
对于$\ell\ge1$, 恒等式$(2\ell+1)P_\ell=P'_{\ell+1}-P'_{\ell-1}$给
$\widetilde A_\ell=Q[P_{\ell+1}(\cos\alpha)-P_{\ell-1}(\cos\alpha)]/[8\pi\eps(2\ell+1)]$.
$\ell=0$另算$\widetilde A_0=Q(1+\cos\alpha)/(8\pi\eps)$, 不用人为矛盾的$P_{-1}$约定. 当$\alpha=0$无孔, 高阶系数全零、恢复均匀球壳; $\alpha=\pi$全挖去、全部系数零. 表面角边缘有理想面荷跳跃, 内部势解不要求边缘处场光滑. 这是按已取得MIT五页讲义重述的例子, 不是凭Jackson题号补造题面.


本章依据 8.07 第5、6讲; 球面边界一般式与作业4完整给出的球腔题对应.
